% TOP_PROBLEM: {"problemId": "problem.hopf-conjecture-for-nonnegative-sectional-curvature", "canonicalTitle": "Hopf conjecture for nonnegative sectional curvature", "releaseRank": 247, "primaryDomain": "geometry_topology", "primaryDomainLabel": "Geometry & topology", "displayStatus": "open", "releaseStatus": "open"}
% !TeX program = lualatex
% Definition and sources only; this file is not an LLM solution attempt.
\documentclass[11pt]{article}
\usepackage[margin=25mm]{geometry}
\usepackage{fontspec}
\setmainfont{Latin Modern Roman}
\usepackage{amsmath,amssymb,mathtools}
\usepackage{unicode-math}
\setmathfont{Latin Modern Math}
\usepackage{xurl,hyperref}
\hypersetup{colorlinks=true,urlcolor=blue}
\setcounter{secnumdepth}{0}
\setlength{\parindent}{0pt}
\setlength{\parskip}{5pt}
\setlength{\emergencystretch}{3em}
\begin{document}
\section*{\texorpdfstring{Hopf conjecture for nonnegative sectional curvature}{Hopf conjecture for nonnegative sectional curvature}}\label{247}
\subsection{Definitions and mathematical statement}
For a closed even-dimensional Riemannian manifold $M$, let $K_g(\sigma)$ be sectional curvature on a tangent two-plane, with the sign convention positive on round spheres, and let
$\chi(M)=\sum_i(-1)^i\dim H_i(M,{\mathbb{Q}})$. The assertion is
\[
 K_g(\sigma)\ge0\text{ on every tangent two-plane}
 \quad\Longrightarrow\quad\chi(M)\ge0.
\]
Zero sectional curvature and zero Euler characteristic are allowed. The manifold need not be simply connected or have an assumed isometric group action. This nonnegative-curvature statement is distinct from the strict positive-curvature version requiring strictly positive Euler characteristic, and neither is merely a scalar-curvature condition.
\par\smallskip\textit{Formulation and concept references:} \href{https://link.springer.com/article/10.1365/s13291-020-00225-x}{[S1]}.

\subsection{Short English statement}
Must every closed even-dimensional manifold with nonnegative sectional curvature have nonnegative Euler characteristic?
\subsection{Sources}
\begin{itemize}
\item[S1]Manuel Amann. Bounds on Sectional Curvature and Interactions with Topology. 2020.
\url{https://link.springer.com/article/10.1365/s13291-020-00225-x}.
\textit{Formulation source.}
\end{itemize}
\end{document}
